% ============================================================================
% SPANISCH LANGENTWURF - DS 6: BRÜCKENTHEMA
% ============================================================================
% Kurs: Jahrgangsübergreifend Spanisch Spätbeginner (Kl. 10, 11, 12)
% Datum: Di 20.01.2026
% Besonderheit: ERSTER OVERLAP zwischen Gruppe A und Gruppe B
% ============================================================================

\documentclass[11pt,a4paper]{article}

% ============================================================================
% PACKAGES
% ============================================================================
\usepackage[utf8]{inputenc}
\usepackage[T1]{fontenc}
\usepackage[ngerman]{babel}
\usepackage{geometry}
\usepackage{fancyhdr}
\usepackage{lastpage}
\usepackage[table]{xcolor}
\usepackage{tcolorbox}
\usepackage{enumitem}
\usepackage{tabularx}
\usepackage{longtable}
\usepackage{array}
\usepackage{booktabs}
\usepackage{hyperref}
\usepackage{ifthen}
\usepackage{amssymb}

% ============================================================================
% KONFIGURATION
% ============================================================================

% --- TIER ---
\newcommand{\plantier}{1}

% --- FACH ---
\newcommand{\fachid}{spanisch}
\newcommand{\fach}{Spanisch}
\newcommand{\fachfarbe}{c41e3a}

% ============================================================================
% VARIABLEN
% ============================================================================

% --- Metadaten ---
\newcommand{\jahrgangsstufe}{10/11/12 (jahrgangsübergreifend)}
\newcommand{\schulart}{Gesamtschule}
\newcommand{\stundenthema}{Brückenthema -- Personas y características}
\newcommand{\reihenthema}{Beobachtungsstunde Februar 2026}
\newcommand{\stundennummer}{6}
\newcommand{\reihenstunden}{8}
\newcommand{\unterrichtsdatum}{20.01.2026}
\newcommand{\unterrichtszeit}{[Lt. Stundenplan]}
\newcommand{\raum}{[Fachraum]}

% --- Lehrkraft ---
\newcommand{\lehrkraft}{[Name]}
\newcommand{\ausbildungsstand}{Lehrkraft}

% --- Lerngruppe ---
\newcommand{\lerngruppengroesse}{10}
\newcommand{\maennlich}{1}
\newcommand{\weiblich}{9}
\newcommand{\divers}{0}

% --- Kompetenzen/Niveau ---
\newcommand{\gerniveau}{A1--A2}
\newcommand{\lehrwerk}{A\_tope.com}
\newcommand{\lehrwerkseite}{S. [X]}

% ============================================================================
% LAYOUT
% ============================================================================
\geometry{left=2.5cm, right=2.5cm, top=2.5cm, bottom=2.5cm}
\definecolor{fach}{HTML}{\fachfarbe}
\definecolor{gray}{HTML}{666666}
\definecolor{lightgray}{HTML}{f5f5f5}
\definecolor{successgreen}{HTML}{2a7a4b}
\definecolor{warningorange}{HTML}{b35c00}
\definecolor{gruppeA}{HTML}{2c5aa0}
\definecolor{gruppeB}{HTML}{c41e3a}
\definecolor{overlap}{HTML}{6a0dad}

\tcbuselibrary{skins,breakable}

% Farbige Boxen
\newtcolorbox{infobox}[1][]{
    colback=lightgray,
    colframe=fach,
    fonttitle=\bfseries,
    title=#1,
    breakable
}

\newtcolorbox{scorebox}{
    colback=white,
    colframe=fach,
    fonttitle=\bfseries,
    title=Didaktik-Score,
    breakable
}

\newtcolorbox{overlapbox}{
    colback=white,
    colframe=overlap,
    fonttitle=\bfseries,
    title={OVERLAP -- Erste gemeinsame Aktivität},
    breakable
}

% Kopf-/Fußzeile
\pagestyle{fancy}
\fancyhf{}
\fancyhead[L]{\small\textcolor{gray}{\fach{} -- Jg. \jahrgangsstufe{} -- \stundenthema}}
\fancyhead[R]{\small\textcolor{gray}{Langentwurf Tier 1}}
\fancyfoot[C]{\small\thepage{} / \pageref{LastPage}}
\renewcommand{\headrulewidth}{0.4pt}
\renewcommand{\footrulewidth}{0pt}

% ============================================================================
% DOKUMENT
% ============================================================================
\begin{document}

% ============================================================================
% DECKBLATT
% ============================================================================
\begin{titlepage}
    \centering
    \vspace*{2cm}

    {\large\textcolor{gray}{\schulart}}\\[2cm]

    {\Huge\textcolor{fach}{\textbf{Langentwurf}}}\\[0.5cm]
    {\large Tier 1 -- Vollständig}\\[2cm]

    {\LARGE\textbf{\stundenthema}}\\[0.5cm]
    {\large im Rahmen der Unterrichtsreihe}\\[0.3cm]
    {\Large\textit{\reihenthema}}\\[1cm]

    \textcolor{overlap}{\Large\textbf{--- ERSTER OVERLAP ---}}\\[0.3cm]
    {\normalsize Erste gemeinsame Aktivität beider Gruppen}\\[1.5cm]

    \begin{tabular}{rl}
        \textbf{Fach:} & \fach{} (\gerniveau) \\[0.3cm]
        \textbf{Klasse:} & \jahrgangsstufe \\[0.3cm]
        \textbf{Lehrwerk:} & \lehrwerk \\[0.3cm]
        \textbf{Datum:} & \underline{\hspace{1cm}\unterrichtsdatum\hspace{1cm}} \\[0.3cm]
        \textbf{Zeit:} & \underline{\hspace{3cm}} (Doppelstunde) \\[0.3cm]
        \textbf{Raum:} & \underline{\hspace{3cm}} \\[0.3cm]
    \end{tabular}

    \vfill

    {\large Lehrkraft: \lehrkraft}\\[0.3cm]
    {\normalsize\textcolor{gray}{\ausbildungsstand}}\\[1cm]

\end{titlepage}

% ============================================================================
% INHALTSVERZEICHNIS
% ============================================================================
\tableofcontents
\newpage

% ============================================================================
% 1. BEDINGUNGSANALYSE
% ============================================================================
\section{Bedingungsanalyse}

\subsection{Lerngruppe}
\begin{tabular}{ll}
    Kursgröße: & \lerngruppengroesse{} SuS (\maennlich{} m / \weiblich{} w / \divers{} d) \\
    Jahrgangsstufe: & Kl. 10, 11, 12 (jahrgangsübergreifend) \\
    Sprachniveau: & \gerniveau{} (GER) -- stark heterogen \\
    Fremdsprache: & 3. Fremdsprache (Spätbeginner) \\
\end{tabular}

\subsection{Schülerprofile und Gruppenzuordnung}

\textbf{Besonderheit dieser Stunde:} Erste gemeinsame Aktivität beider Gruppen (OVERLAP)!

\vspace{0.5em}

\begin{tabular}{|l|c|l|p{5.5cm}|}
\hline
\textbf{Name} & \textbf{Kl.} & \textbf{Gruppe} & \textbf{Profil} \\
\hline
\textcolor{gruppeA}{E.H.} & 10 & \textcolor{gruppeA}{A} & TOP, 15 Punkte, Franz. Hintergrund \\
\hline
\textcolor{gruppeA}{L.N.} & 10 & \textcolor{gruppeA}{A} & Proaktiv, etwas faul \\
\hline
\textcolor{gruppeA}{C.P.} & 10 & \textcolor{gruppeA}{A} & 12 Punkte, ruhig, gute Aussprache \\
\hline
\textcolor{gruppeA}{V.S.} & 10 & \textcolor{gruppeA}{A} & Konstant, fleißig \\
\hline
\textcolor{gruppeB}{L.K.} & 11 & \textcolor{gruppeB}{B} & 10 Punkte, solide \\
\hline
\textcolor{gruppeB}{H.P.} & 11 & \textcolor{gruppeB}{B} & Nachzügler, 10 Punkte \\
\hline
\textcolor{gruppeB}{A.A.} & 12 & \textcolor{gruppeB}{B} & Fleißig, unsicher \\
\hline
\textcolor{gruppeB}{A.S.} & 12 & \textcolor{gruppeB}{B} & 9 Punkte, einziger Junge \\
\hline
\end{tabular}

\vspace{1em}

\textbf{Gruppen-Charakteristik:}
\begin{itemize}
    \item \textcolor{gruppeA}{\textbf{Gruppe A (4 SuS, Kl. 10):}} Starke Gruppe, homogenes Niveau (A2), gute Aussprache, können frei sprechen
    \item \textcolor{gruppeB}{\textbf{Gruppe B (4 SuS, Kl. 11/12):}} Heterogen, A1--A2, teilweise unsicher, benötigen mehr Struktur
\end{itemize}

\subsection{Differenzierung (Brückenthema)}

\textbf{Ziel des Brückenthemas:} Gemeinsame sprachliche Basis schaffen für OVERLAP

\begin{tabular}{|l|p{5cm}|p{5cm}|}
\hline
& \textcolor{gruppeA}{\textbf{Gruppe A}} & \textcolor{gruppeB}{\textbf{Gruppe B}} \\
\hline
\textbf{Fokus} & Personen beschreiben (Familie, Freunde) & Eigenschaften für Berufe \\
\hline
\textbf{Niveau} & Freie Produktion mit Konnektoren & Strukturierte Satzbausteine \\
\hline
\textbf{Ziel} & Flüssig über Menschen sprechen & Adjektive + hay que ser + porque... \\
\hline
\end{tabular}

\subsection{Besonderheiten}
\begin{itemize}
    \item \textbf{Heterogenität:} 4 Lernjahre Unterschied innerhalb des Kurses
    \item \textbf{Overlap-Vorbereitung:} Erste gemeinsame Aktivität muss niedrigschwellig sein
    \item \textbf{Sozial:} Gruppen kennen sich kaum -- Eisbrecher-Funktion wichtig
    \item \textbf{A.S.:} Einziger Junge, Gruppenkonstellationen beachten
\end{itemize}

% ============================================================================
% 2. SACHANALYSE
% ============================================================================
\section{Sachanalyse}

\subsection{Sprachliche Strukturen}

\textbf{Kernthema:} Personenbeschreibung mit Adjektiven

\begin{itemize}
    \item \textbf{Grammatik:}
    \begin{itemize}
        \item \textit{ser + Adjektiv} (Eigenschaften): \textit{Mi hermana es simpática.}
        \item \textit{estar + Adjektiv} (Zustände): \textit{Hoy estoy cansado.}
        \item \textit{hay que + Infinitiv} (Notwendigkeit): \textit{Hay que ser paciente.}
        \item \textit{porque} (Begründung): \textit{...porque trabaja con niños.}
    \end{itemize}

    \item \textbf{Wortschatz (Adjektive):}
    \begin{itemize}
        \item Positiv: \textit{simpático/-a, inteligente, trabajador/-a, responsable, paciente, amable, divertido/-a, creativo/-a}
        \item Neutral: \textit{tranquilo/-a, serio/-a, reservado/-a, organizado/-a}
        \item Kritisch: \textit{perezoso/-a, impaciente, egoísta}
    \end{itemize}

    \item \textbf{Konnektoren:}
    \begin{itemize}
        \item \textit{también} (auch), \textit{además} (außerdem), \textit{pero} (aber)
        \item \textit{porque} (weil), \textit{por eso} (deshalb)
    \end{itemize}
\end{itemize}

\subsection{Kontrastive Betrachtung (Deutsch $\leftrightarrow$ Spanisch)}
\begin{itemize}
    \item \textbf{Positive Transfers:}
    \begin{itemize}
        \item Viele Adjektive sind international verständlich (\textit{inteligente, responsable})
        \item Satzstruktur bei \textit{porque} ähnlich wie bei ``weil''
    \end{itemize}

    \item \textbf{Interferenzen:}
    \begin{itemize}
        \item \textit{ser} vs. \textit{estar}: Im Deutschen nur ``sein''
        \item Adjektivendungen (-o/-a) fehlen im Deutschen
        \item \textit{hay que} hat keine direkte deutsche Entsprechung
    \end{itemize}

    \item \textbf{Häufige Fehler:}
    \begin{itemize}
        \item *\textit{Mi hermano es simpático-a} (Gender-Endung vergessen/falsch)
        \item *\textit{Estoy inteligente} (\textit{ser}/\textit{estar}-Verwechslung)
        \item *\textit{Hay que es paciente} (falscher Anschluss)
    \end{itemize}
\end{itemize}

\subsection{Kulturelle Aspekte}
\begin{itemize}
    \item Persönliche Fragen nach Familie sind in spanischsprachigen Kulturen üblich
    \item Offenheit bei Personenbeschreibungen (auch körperliche Merkmale)
    \item Berufe und soziale Rollen: Traditionelle vs. moderne Erwartungen
\end{itemize}

% ============================================================================
% 3. DIDAKTISCHE ANALYSE
% ============================================================================
\section{Didaktische Analyse}

\subsection{Stellung in der Sequenz}

Dies ist die \textbf{6. Doppelstunde} der Beobachtungssequenz (8 DS gesamt).

\vspace{0.5em}
\textbf{Sequenzübersicht:}
\begin{center}
\begin{tabular}{|c|l|l|}
\hline
\textbf{DS} & \textbf{Thema} & \textbf{Fokus} \\
\hline
1--5 & Getrennte Progression & A: LJ1 / B: LJ2 \\
\hline
\rowcolor{lightgray} \textbf{6} & \textbf{Brückenthema: Personas y características} & \textbf{ERSTER OVERLAP} \\
\hline
7 & Beobachtungsstunde & Vollständige Integration \\
\hline
8 & Abschluss & Reflexion, Evaluation \\
\hline
\end{tabular}
\end{center}

\textbf{Funktion dieser Stunde:}
\begin{itemize}
    \item \textbf{Vorbereitung:} Brückenthema schafft gemeinsamen Wortschatz
    \item \textbf{Erprobung:} Erster Overlap testet Gruppendynamik
    \item \textbf{Diagnostik:} Beobachtung der Interaktion für DS 7
\end{itemize}

\subsection{Kompetenzziele (Rahmenplan MV)}
\begin{itemize}
    \item \textbf{Sprechen (dialogisch):} Über Personen und ihre Eigenschaften sprechen (A1.2/A2.1)
    \item \textbf{Hörverstehen:} Einfache Personenbeschreibungen verstehen
    \item \textbf{Sprachliche Mittel:} \textit{ser} + Adjektiv, \textit{hay que} + Infinitiv
    \item \textbf{Interkulturell:} Offenheit für verschiedene Perspektiven auf Eigenschaften
    \item \textbf{Methodisch:} Partnerinterview, Notizen anfertigen
\end{itemize}

\subsection{Legitimation}

\textbf{Rahmenplan Spanisch MV (Gymnasiale Oberstufe):}
\begin{itemize}
    \item Kompetenzbereich: Kommunikative Fertigkeiten -- Sprechen (dialogisch)
    \item Standard: ``Einfache Informationen über Personen austauschen''
    \item Sprachliche Mittel: Adjektive zur Personenbeschreibung
\end{itemize}

\textbf{Brückenthemen-Legitimation:}
Das gemeinsame Thema ``Personas y características'' ermöglicht:
\begin{itemize}
    \item Anknüpfung an Vorwissen beider Gruppen
    \item Niedrigschwelligen Einstieg für Gruppe B (Fragen stellen)
    \item Anspruchsvolle Erweiterung für Gruppe A (freie Antworten)
\end{itemize}

% ============================================================================
% 4. STUNDENZIELE
% ============================================================================
\section{Stundenziele}

\subsection{Grobziel}
Die SuS erweitern ihre \textbf{dialogische Sprechkompetenz}, indem sie Personen mit passenden Adjektiven beschreiben und in einem gruppenübergreifenden Interview anwenden.

\subsection{Feinziele (operationalisiert)}

\begin{tabular}{|c|p{7cm}|c|c|c|}
\hline
\textbf{Nr.} & \textbf{Die SuS können...} & \textbf{AFB} & \textbf{Gruppe} & \textbf{Phase} \\
\hline
FZ1 & mindestens 8 Adjektive zur Personenbeschreibung korrekt verwenden & I & alle & Einstieg \\
\hline
FZ2 & den Unterschied \textit{ser/estar} + Adj. an Beispielen erklären & II & alle & Einstieg \\
\hline
FZ3 & eine Person mit mind. 3 Adjektiven + Konnektoren beschreiben & II & A & Erarbeitung \\
\hline
FZ4 & berufliche Eigenschaften mit \textit{hay que ser... porque...} formulieren & II & B & Erarbeitung \\
\hline
FZ5 & ein einfaches Interview auf Spanisch führen und Notizen anfertigen & II & alle & \textcolor{overlap}{Overlap} \\
\hline
FZ6 & eine schriftliche Zusammenfassung über einen Mitschüler verfassen & III & B & \textcolor{overlap}{Overlap} \\
\hline
\end{tabular}

% ============================================================================
% 5. METHODISCHE ANALYSE
% ============================================================================
\section{Methodische Analyse}

\subsection{Phasierung (Doppelstunde = 90 Min.)}

\begin{center}
\begin{tabular}{|c|l|c|l|}
\hline
\textbf{Zeit} & \textbf{Phase} & \textbf{Min.} & \textbf{Sozialform} \\
\hline
0--15 & Gemeinsamer Einstieg & 15 & Plenum \\
\hline
15--45 & Getrennte Erarbeitung (A + B parallel) & 30 & GA/PA \\
\hline
45--50 & Kurze Pause & 5 & -- \\
\hline
50--55 & Overlap-Vorbereitung & 5 & Plenum \\
\hline
\rowcolor{lightgray} 55--70 & \textcolor{overlap}{\textbf{OVERLAP: Experten-Interview}} & 15 & \textcolor{overlap}{Tandem (A+B)} \\
\hline
70--85 & Präsentation + Feedback & 15 & Plenum \\
\hline
85--90 & Abschluss + HA & 5 & Plenum \\
\hline
\end{tabular}
\end{center}

\subsection{Detailplanung Overlap (15 Min.)}

\begin{overlapbox}
\textbf{Aktivität: ``Expertos entrevistan'' (Experten-Interview)}

\vspace{0.5em}
\textbf{Ablauf:}
\begin{enumerate}
    \item B-Schüler interviewen A-Schüler (1:1 oder 2:1)
    \item B stellt einfache Fragen: \textit{¿Cómo eres? ¿Cómo es tu familia?}
    \item A antwortet frei mit Adjektiven + Konnektoren
    \item B notiert Stichpunkte und schreibt anschließend Zusammenfassung
\end{enumerate}

\textbf{Paarzuordnung (Vorschlag):}
\begin{tabular}{ll}
\textcolor{gruppeB}{L.K.} $\rightarrow$ \textcolor{gruppeA}{E.H.} & (solide $\leftarrow$ TOP) \\
\textcolor{gruppeB}{H.P.} $\rightarrow$ \textcolor{gruppeA}{C.P.} & (Nachzügler $\leftarrow$ ruhig, geduldig) \\
\textcolor{gruppeB}{A.A.} $\rightarrow$ \textcolor{gruppeA}{V.S.} & (unsicher $\leftarrow$ fleißig, verständnisvoll) \\
\textcolor{gruppeB}{A.S.} $\rightarrow$ \textcolor{gruppeA}{L.N.} & (Junge $\leftarrow$ proaktiv) \\
\end{tabular}

\vspace{0.5em}
\textbf{Scaffolding für Gruppe B:}
\begin{itemize}
    \item Fragekarten: \textit{¿Cómo eres? / ¿Y tu mejor amigo? / ¿Por qué?}
    \item Notizvorlage: ``[Nombre] es \_\_\_ y \_\_\_. También es \_\_\_.''
\end{itemize}

\textbf{Ziel:} Erste erfolgreiche Zusammenarbeit $\rightarrow$ Vorbereitung auf DS 7 (Beobachtungsstunde)
\end{overlapbox}

\subsection{Medien}
\begin{itemize}
    \item Tafel/Whiteboard: Adjektiv-Mindmap
    \item Wortkarten: Adjektive (laminiert, zum Sortieren)
    \item AB-06a: Personenbeschreibung (Gruppe A)
    \item AB-06b: Berufe + Eigenschaften (Gruppe B)
    \item Interview-Fragekarten (Gruppe B)
    \item Notizvorlage für Overlap
\end{itemize}

% ============================================================================
% 6. VERLAUFSPLANUNG
% ============================================================================
\section{Verlaufsplanung}

\textbf{Übersicht Doppelstunde (90 Min.)}

\begin{longtable}{|p{0.8cm}|p{1.5cm}|p{4cm}|p{3cm}|p{1.5cm}|p{2cm}|}
\hline
\textbf{Zeit} & \textbf{Phase} & \textbf{Lehrerhandeln} & \textbf{Schülerhandeln} & \textbf{SF} & \textbf{Medien} \\
\hline
\endhead

\multicolumn{6}{|c|}{\textbf{GEMEINSAMER EINSTIEG (15 Min.)}} \\
\hline

0--3 & Einstieg & Begrüßung: ``¡Buenos días! Hoy trabajamos juntos.'' Ankündigung Brückenthema & Zuhören, sich sammeln & Plenum & -- \\
\hline

3--8 & Aktivierung & Impuls: ``¿Cómo es una persona simpática?'' Adjektive sammeln an der Tafel & Nennen Adjektive, L notiert & Plenum & Tafel \\
\hline

8--13 & Input & Unterschied \textit{ser} + Adj. vs. \textit{estar} + Adj. klären mit Beispielen & Beispiele nachvollziehen, eigene Beispiele nennen & Plenum & Tafel \\
\hline

13--15 & Überleitung & Gruppenaufteilung erklären: ``Grupo A: personas. Grupo B: profesiones.'' & In Gruppen aufteilen & Plenum & -- \\
\hline

\multicolumn{6}{|c|}{\textcolor{gruppeA}{\textbf{GRUPPE A: Personenbeschreibung (30 Min., parallel zu B)}}} \\
\hline

15--22 & Erarbeitung & \textcolor{gruppeA}{[A] Impuls: ``Describe a tu hermano/amigo''}, Wortschatz erweitern & \textcolor{gruppeA}{Notieren Beschreibungen} & EA & AB-06a \\
\hline

22--30 & Übung 1 & \textcolor{gruppeA}{Partner-Interview anleiten: ``¿Cómo es tu mejor amigo/-a?''} & \textcolor{gruppeA}{Gegenseitig interviewen} & PA & -- \\
\hline

30--40 & Vertiefung & \textcolor{gruppeA}{Konnektoren einführen: \textit{también, además, pero}} & \textcolor{gruppeA}{Erweiterte Beschreibungen üben} & PA & Tafel \\
\hline

40--45 & Sicherung & \textcolor{gruppeA}{2--3 SuS stellen Partner vor} & \textcolor{gruppeA}{Präsentieren} & Plenum & -- \\
\hline

\multicolumn{6}{|c|}{\textcolor{gruppeB}{\textbf{GRUPPE B: Eigenschaften für Berufe (30 Min., parallel zu A)}}} \\
\hline

15--22 & Erarbeitung & \textcolor{gruppeB}{[B] Berufe nennen, Eigenschaften zuordnen} & \textcolor{gruppeB}{Brainstorming, Zuordnung} & Plenum B & Wortkarten \\
\hline

22--32 & Input & \textcolor{gruppeB}{\textit{Para ser médico hay que ser paciente porque...}} & \textcolor{gruppeB}{Struktursätze bilden} & EA/PA & AB-06b \\
\hline

32--40 & Übung & \textcolor{gruppeB}{Mini-Diskussion: ``¿Un buen jefe?''} & \textcolor{gruppeB}{Diskutieren in kleiner Runde} & GA & -- \\
\hline

40--45 & Sicherung & \textcolor{gruppeB}{Zusammenfassung, beste Sätze notieren} & \textcolor{gruppeB}{Notieren} & Plenum B & Tafel \\
\hline

\multicolumn{6}{|c|}{\textbf{PAUSE (5 Min.)}} \\
\hline

45--50 & Pause & Pause, Raum vorbereiten für Overlap & Pause & -- & -- \\
\hline

\multicolumn{6}{|c|}{\textbf{OVERLAP-VORBEREITUNG (5 Min.)}} \\
\hline

50--55 & Vorbereitung & Erklärung Experten-Interview: B interviewt A. Paare zuordnen, Fragekarten verteilen & Paare finden, Material sichten & Plenum & Fragekarten \\
\hline

\multicolumn{6}{|c|}{\textcolor{overlap}{\textbf{OVERLAP: EXPERTEN-INTERVIEW (15 Min.)}}} \\
\hline

55--65 & \textcolor{overlap}{Interview} & \textcolor{overlap}{Beobachten, bei Bedarf helfen, Notizen für Feedback} & \textcolor{overlap}{B interviewt A auf Spanisch, notiert Stichpunkte} & \textcolor{overlap}{Tandem} & \textcolor{overlap}{Fragekarten, Notizvorlage} \\
\hline

65--70 & \textcolor{overlap}{Schreiben} & \textcolor{overlap}{B-SuS schreiben Zusammenfassung, A-SuS helfen bei Vokabeln} & \textcolor{overlap}{B schreibt: ``[Name] es...''} & \textcolor{overlap}{Tandem} & \textcolor{overlap}{Notizvorlage} \\
\hline

\multicolumn{6}{|c|}{\textbf{PRÄSENTATION + FEEDBACK (15 Min.)}} \\
\hline

70--80 & Präsentation & Ausgewählte B-SuS lesen Zusammenfassung vor, A-SuS bestätigen/ergänzen & 3--4 Tandems präsentieren & Plenum & -- \\
\hline

80--85 & Feedback & Positives Feedback, sprachliche Korrektur (sanft), Reflexion Zusammenarbeit & Zuhören, reflektieren & Plenum & -- \\
\hline

\multicolumn{6}{|c|}{\textbf{ABSCHLUSS (5 Min.)}} \\
\hline

85--88 & Ausblick & ``Nächste Woche: Beobachtungsstunde mit noch mehr Zusammenarbeit!'' & Zuhören & Plenum & -- \\
\hline

88--90 & HA & HA: 3 Sätze über Familienmitglied (A) / 3 Sätze über Traumberuf (B) & Notieren & Plenum & -- \\
\hline

\end{longtable}

\textbf{Legende:} EA = Einzelarbeit, PA = Partnerarbeit, GA = Gruppenarbeit, SF = Sozialform

% ============================================================================
% 7. ERWARTETE SCHWIERIGKEITEN
% ============================================================================
\section{Erwartete Schwierigkeiten}

\begin{tabular}{|p{4.5cm}|p{8.5cm}|}
\hline
\textbf{Schwierigkeit} & \textbf{Reaktion} \\
\hline
\textbf{B-SuS unsicher beim Fragen stellen} & Fragekarten mit vorformulierten Fragen bereitstellen. L modelliert Interview mit einem A-Schüler vor allen. \\
\hline
\textbf{A-SuS antworten zu schnell/komplex} & A-SuS explizit bitten, langsam und deutlich zu sprechen. Vorab Anweisung: ``Sprich so, dass B dich versteht.'' \\
\hline
\textbf{Paarfindung problematisch} & Feste Zuordnung durch L (siehe Vorschlag in 5.2). A.S. als einziger Junge gezielt zu L.N. (proaktiv, unkompliziert). \\
\hline
\textbf{Zeitproblem bei Overlap} & Puffer: Präsentation auf 2--3 Tandems reduzieren statt 4. Schreibphase in HA auslagern, falls nötig. \\
\hline
\textbf{\textit{ser/estar}-Verwechslung} & Merkhilfe: ``\textit{Ser} = so bin ich IMMER, \textit{estar} = so bin ich JETZT.'' Tafelanschrieb sichtbar lassen. \\
\hline
\textbf{Adjektiv-Endungen (-o/-a)} & Bei Korrektur auf Gender des Subjekts hinweisen. Visuelle Hilfe: Adjektive mit -o (blau) / -a (rot) markieren. \\
\hline
\textbf{B-SuS verstehen A-Antworten nicht} & A-SuS dürfen paraphrasieren oder einzelne Wörter auf Deutsch erklären. L als Backup verfügbar. \\
\hline
\end{tabular}

% ============================================================================
% 8. HAUSAUFGABE
% ============================================================================
\section{Hausaufgabe}

\textbf{Differenziert nach Gruppen:}

\begin{tabular}{|l|p{10cm}|}
\hline
\textcolor{gruppeA}{\textbf{Gruppe A}} & Schreibe 3 Sätze über ein Familienmitglied mit mindestens 3 verschiedenen Adjektiven und Konnektoren. \\
& \textit{Beispiel: Mi madre es muy trabajadora. También es paciente, pero a veces es impaciente por las mañanas.} \\
\hline
\textcolor{gruppeB}{\textbf{Gruppe B}} & Schreibe 3 Sätze über deinen Traumberuf mit \textit{hay que ser... porque...} \\
& \textit{Beispiel: Para ser profesor hay que ser paciente porque trabajas con muchos niños.} \\
\hline
\end{tabular}

\vspace{1em}
\textbf{Alle:} Vokabeln wiederholen: Adjektive zur Personenbeschreibung

% ============================================================================
% 9. ANLAGEN
% ============================================================================
\section{Anlagen}

\begin{enumerate}
    \item AB-06a: Personenbeschreibung (Gruppe A)
    \item AB-06b: Berufe + Eigenschaften (Gruppe B)
    \item Interview-Fragekarten (für Gruppe B)
    \item Notizvorlage Overlap (für Gruppe B)
    \item Wortkarten Adjektive (laminiert)
    \item Tafelbild: \textit{ser} vs. \textit{estar} + Adjektiv
\end{enumerate}

\subsection{Tafelbild (Entwurf)}

\begin{verbatim}
+---------------------------------------------------------------------+
|                    ¿Cómo es una persona?                            |
+---------------------------------------------------------------------+
|                                                                     |
|  SER + Adjektiv (= Eigenschaft)    |   ESTAR + Adjektiv (= Zustand) |
|  ------------------------------    |   ----------------------------  |
|  Soy simpático. (Ich bin nett.)    |   Estoy cansado. (Ich bin müde.)|
|  Es inteligente. (Er ist klug.)    |   Está enfermo. (Er ist krank.) |
|  = So bin ich IMMER                |   = So bin ich JETZT            |
|                                                                     |
+---------------------------------------------------------------------+
|  Adjektive:                                                         |
|  simpático/a - inteligente - trabajador/a - paciente - amable       |
|  divertido/a - creativo/a - responsable - organizado/a - serio/a    |
|                                                                     |
|  Konnektoren: también (auch) - además (außerdem) - pero (aber)      |
|                                                                     |
|  Struktur B:  Para ser [Beruf] hay que ser [Adj.] porque...         |
+---------------------------------------------------------------------+
\end{verbatim}

\subsection{Interview-Fragekarten (für Gruppe B)}

\begin{verbatim}
+------------------------------------+
|  ¿Cómo eres?                       |
|  (Wie bist du?)                    |
+------------------------------------+

+------------------------------------+
|  ¿Cómo es tu mejor amigo/-a?       |
|  (Wie ist dein bester Freund?)     |
+------------------------------------+

+------------------------------------+
|  ¿Cómo es tu familia?              |
|  (Wie ist deine Familie?)          |
+------------------------------------+

+------------------------------------+
|  ¿Por qué?                         |
|  (Warum?)                          |
+------------------------------------+
\end{verbatim}

\subsection{Notizvorlage Overlap}

\begin{verbatim}
+---------------------------------------------------------------------+
|  NOTIZEN: Interview mit _______________________                      |
+---------------------------------------------------------------------+
|                                                                     |
|  Er/Sie ist: ________________________________________________       |
|                                                                     |
|  Sein/Ihr bester Freund ist: ____________________________________   |
|                                                                     |
|  Seine/Ihre Familie ist: ________________________________________   |
|                                                                     |
+---------------------------------------------------------------------+
|  ZUSAMMENFASSUNG (auf Spanisch):                                    |
|                                                                     |
|  _____________ es _____________ y _____________.                    |
|                                                                     |
|  También es _____________.                                          |
|                                                                     |
|  Su mejor amigo/-a es _____________.                                |
|                                                                     |
+---------------------------------------------------------------------+
\end{verbatim}

% ============================================================================
% 10. DIDAKTIK-SCORE
% ============================================================================
\newpage
\section{Didaktik-Score}

\begin{scorebox}
\textbf{Bewertung der didaktischen Qualität nach 10 Dimensionen}\\
\textbf{Ziel-Score: $\geq$ 9.0 (Premium-Qualität)}

\vspace{1em}
\begin{tabular}{|c|l|c|c|p{4.5cm}|}
\hline
\textbf{\#} & \textbf{Dimension} & \textbf{Gew.} & \textbf{Score} & \textbf{Notiz} \\
\hline
1 & Differenzierung & 15\% & 10/10 & A/B parallel + Overlap mit gestuften Rollen \\
\hline
2 & Taxonomie (AFB) & 10\% & 9/10 & I-III abgedeckt, AFB III in Overlap \\
\hline
3 & Lernziel-Alignment & 10\% & 10/10 & Passt zu RP MV, A1-A2 Niveau \\
\hline
4 & Aktivierung & 10\% & 10/10 & Interview = konstruktiv/interaktiv \\
\hline
5 & Scaffolding & 10\% & 10/10 & Fragekarten, Notizvorlage, L-Modellierung \\
\hline
6 & Feedback-Qualität & 10\% & 9/10 & Peer-Feedback + L-Korrektur, elaboriert \\
\hline
7 & Sprachliche Klarheit & 10\% & 10/10 & Altersgerecht, klare Strukturen \\
\hline
8 & Motivation \& Relevanz & 10\% & 9/10 & Persönlicher Bezug (Familie, Freunde) \\
\hline
9 & Inklusion (UDL) & 10\% & 9/10 & Multiple Zugänge, aber auditiv dominant \\
\hline
10 & Praktikabilität & 5\% & 10/10 & Zeitrahmen realistisch, Material vorhanden \\
\hline
\multicolumn{3}{|r|}{\textbf{GESAMT:}} & \textbf{9.6/10} & \\
\hline
\end{tabular}

\vspace{1em}
\textbf{Bewertungsschwellen:}
\begin{itemize}
    \item[\checkmark] \textcolor{successgreen}{$\geq$ 9.0: \textbf{FREIGABE} (Premium-Qualität)}
    \item[$\Box$] \textcolor{warningorange}{8.0--8.9: Iteration empfohlen}
    \item[$\Box$] 6.0--7.9: Überarbeitung erforderlich
    \item[$\Box$] $<$ 6.0: Grundlegende Neukonzeption
\end{itemize}

\vspace{1em}
\textcolor{successgreen}{\textbf{STATUS: FREIGABE}}

\vspace{1em}
\textbf{Stärken:}
\begin{itemize}
    \item Klare Brückenthema-Konzeption mit echtem OVERLAP
    \item Sorgfältige Paarzuordnung basierend auf Schülerprofilen
    \item Umfassendes Scaffolding für unsichere B-Schüler
    \item Vorbereitung auf Beobachtungsstunde (DS 7) explizit
\end{itemize}

\textbf{Optionale Verbesserungen:}
\begin{itemize}
    \item Visuelle Unterstützung (Bilder von Berufen) ergänzen
    \item Audio-Beispiel für \textit{ser/estar}-Kontrast
\end{itemize}

\end{scorebox}

\end{document}
